\chapter{实现与交互演示}
\label{ch:demo}

正文中的轨迹、交叉表和命题核对来自同一套数值核。交互页面读取该核已经写出的状态流，在三维腔体里同时显示四个边界、食品温度色带、因果路径的显示量和安全投影前后的加热。着色器颜色是两个标量的插值，不是现场积分。把交叉表读回可见零件之后，加热管发光、蒸汽从喷嘴喷出、风机叶片加速、排汽百叶张开，底部曲线同时给出表面、冷点、估计和进度。数字与附录\ref{app:trace} 的抽样一致。

\section{数值核}

核沿厚度用十六条控制体，时间步 \(2\,\mathrm{s}\)，时域 55 min。状态包括温度、含水、加热器温度、腔温腔湿、热阻代理、冷点跟踪器和进度。四种策略与三档允许集合写在同一循环里。噪声是确定性散列，重复同一种子得到同一轨迹。离散方程见附录\ref{app:discrete}。

数值核把状态写成 JSON 与 CSV。演示读取 \texttt{catalog.json}，按策略、档位和厚度选择已经写出的数组，帧间隔 8 s，与附录\ref{app:trace} 的抽样一致。默认条件是组合档、对偶策略、肉糜 30 mm、初温 \(8^\circ\mathrm{C}\)、含水 0.72，对应种子 7 那一条已经写进附录的轨迹。catalog 预写 12 条：四种策略乘三种档位给出 30 mm 默认，再加上 15 mm 与 50 mm 的对偶组合。厚度滑条只切换所读数组和几何缩放。

命题核对是独立测试：平衡观测对内部热阻的极差为 0；两段脉冲的内部滞后为 40.846 与 38.912，差值大于 1；损伤上界 1.2 盖住离散实现 1.143；屏障测试函数不随后验改变。测试几何见各命题的证明，常数不必等于 S1 的 40 s 窗口。

\section{三维腔体}

腔体按家用蒸烤箱的可见零件建模。外壳为深色钢板，夹层为岩棉色保温块，内衬为圆角拉丝不锈钢，底部另有一层台面衬板，用来把腔体从地面抬开。门框用黄铜条和橡胶密封条勾边，门玻璃做成内外两层，中间留一条窄气隙，铰链侧打开约 \(0.72\,\mathrm{rad}\)，把手为镀铬圆杆。顶板前方四档旋钮对应加热、蒸汽、风机、模式，控制窗标出 ACTS-MCDT 2402。顶部六根 U 形加热管带陶瓷端帽、镀铬插针和前后汇流排，管壁自发光强度跟加热指令走。后壁是十叶扭曲翼型风机：叶片用深灰金属色、更厚的挤压截面，放在四辐护网前方，转速跟风机指令走；护网后面另有一层近黑挡板和电机罩，把叶片从拉丝内衬上分开，默认摄像机方位下叶片轮廓可辨。左侧蒸汽管通向双锥喷嘴；蒸汽指令大于 0.5 时，喷嘴锥与约 720 个加性粒子同时可见，基线蒸汽 0.15 只保留管路、不喷雾。右上七片排汽百叶随排汽指令张开。腔内有五档搁架、带冲孔的烤盘、穿刺探针、水箱液面（随水位缩放）、冷凝托盘、腔灯、铰链侧限位和四只底脚。色调映射为 ACES，阴影为 PCF 软阴影，环境用 PMREM。

零件数量与正文图\ref{fig:cavity} 的四个执行量一一对应：六管是加热，喷嘴与粒子是蒸汽，十叶是风机，七片百叶是排汽。水位只出现在水箱液面，门只出现在铰链开度，二者进入安全侧测量。摄像机可以绕腔体旋转，用来核对加热管是否挡住食品色带、蒸汽粒子是否从左侧喷嘴出发。默认方位角 0.52、俯角 0.20、距离 6.55。

食品沿厚度分成十六块几乎贴合的圆角薄片，与控制体一一对应。相邻薄片的间隙压到控制体宽度的百分之一量级，十六块连成一块带色带的整肉。第 \(i\) 片的颜色由
\[
T_i=T_s-(T_s-T_{\mathrm{cold}})\,u_i^{1.15},\qquad u_i=i/15
\]
插值，\(T_s\) 与 \(T_{\mathrm{cold}}\) 来自当前帧日志。第 5 片对应脂肪夹层，用更深的色带标出。厚度滑条只改变几何缩放和所读数组。横向没有第二维梯度。加热指令提高管壁自发光、腔灯和补光；风机指令提高叶片转速；排汽指令拉伸百叶。这些零件把交叉表里的四个执行量读回可见几何。

\section{面板、对照与事件}

页面分成三栏。左栏选择蒸、烤、组合与四种策略，并给出时间、厚度、对照策略和门开度。中栏是腔体。右栏给出本帧的含水、进度、热阻与方差、腔温腔湿、水位，以及四个执行器的条形指示。路径 A/B/C 的显示量来自该帧已经写出的通量代理。请求与指令分两列，用来看见安全投影有没有改写加热；本次工作点上二者重合，因为加热器停在 \(230^\circ\mathrm{C}\) 以下。

底部曲线同时画出表面温度、冷点温度、估计温度和进度。左轴为 \(0\)--\(180^\circ\mathrm{C}\)，右轴为进度 \(\alpha\in[0,1]\)，横轴为分钟；浅色竖带标出 150--190 s 探针窗。若打开对照策略，第二条冷点与进度以虚线叠在同一坐标上，用来直接看见开环菜单切断、腔温比例不停机、无探针提前停、对偶窗后方差下降这四件事。事件条标出探针窗、1680 s 品质段、停机和保持；跳转按钮把滑条送到这些时刻。门开关只改铰链开度，进入安全侧测量，不改已经写出的温度数组。

回放间隔 40 ms，整条 55 min 轨迹约半分钟走完。探针窗内芯片切到“探针”，蒸汽条拉满；窗后芯片回到“品质”，加热请求由估计温度决定。切入保持之后加热降到 0.15，色带从表面向冷点回落。这一顺序与第\ref{ch:ati}章的时间表一致。

读演示时按三帧走即可核对附录。起始帧：加热请求与指令重合为 1，蒸汽为基线 0.15，十六块薄片接近初温，进度为 0。探针窗第 20 帧（约 160 s）：蒸汽粒子从左侧喷嘴喷出，表面已经显著高于冷点，估计仍接近零。晚期第 210 帧（1680 s）：加热仍在品质段，冷点 \(70.21^\circ\mathrm{C}\)，进度 0.558，蒸汽回到基线。左栏同时给出本条轨迹的停机时刻、\(e_{\alpha}\) 和探针是否把方差降到 0.224。导出 CSV 的字段与核日志一致，并附带 \texttt{summary.json} 与 \texttt{parameters.json}。

图\ref{fig:demo}给出组合档、对偶策略、30 mm 的三帧。厚度切到 50 mm 时，同一策略在时域终点进度只有 0.068，色带冷点一侧几乎仍停在初温附近：探针已经把方差降到 0.224，导热时间尺度没有走完，见图\ref{fig:demothick}。对照开环时，1170 s 加热切断，冷点停在 \(16^\circ\mathrm{C}\) 量级。停机帧把进度送到上沿附近：页面按 8 s 抽样取切入保持前一帧，显示 1864 s、\(e_{\alpha}=+0.145\)，核日志的停机是 1866 s，见图\ref{fig:demostop}。15 mm / 30 mm / 50 mm 三档芯片只切换已经写出的数组。这三次切换把第\ref{ch:result}章的厚度符号和第\ref{ch:atlas}章的策略轨迹读回同一腔体。

\section{交互页与导出字段}

交互页读取已经写出的 \texttt{catalog.json}，不在浏览器里重跑十六条积分。对照、跳转、厚度三档和开门是对同一数组的读法，不是第二条求解器。公开地址为
\begin{center}
\url{https://2711944586.github.io/xxgz/}
\end{center}
页面顶栏重复写出该地址。默认条件对应种子 7、肉糜 30 mm、组合档、对偶策略；厚度滑条只切换所读数组和几何缩放。

页面导出二十二列，顺序是时间、表面温度、冷点、估计、含水、进度、热阻代理、方差、加热指令、加热请求、蒸汽、风机、排汽、模式、加热器温度、水位、腔温、腔湿、壳体、路径 A/B/C。字段名写成一行就是
\begin{quote}\small\ttfamily\raggedright
time, t1, t2, hat, m, alpha, rHat, varR, heat, heatReq, steam, fan, exhaust, mode, heater, water, cavT, cavH, shell, pathA, pathB, pathC.
\end{quote}
步长 \(2\,\mathrm{s}\)，每四步记一行，间隔 \(8\,\mathrm{s}\)，默认时长 55 min，与种子 7 的 Python 轨迹相同，共 414 行。\(\texttt{t2}\) 是十六条里更新前的最低温度，\(\texttt{hat}\) 才是估计。两列都在，不能把冷点列当成估计。Python 的 \texttt{demo\_seed7\_dual.csv} 有平均含水，没有加热器。两边共有的表面、冷点和估计，在种子 7、肉糜、30 mm、初温 \(8^\circ\mathrm{C}\)、含水 0.72、对偶策略上，最大差低于 \(5\times 10^{-6}\,\mathrm{K}\)。差来自 CSV 保留五位小数。浏览器上改过厚度或策略之后的导出，没有再和 Python 对过。能量若从浏览器的加热列重算，得到的是指令乘 \(8\,\mathrm{s}\) 的展示和。规格中的 \(e_\alpha\) 仍是停止时刻的 \(\alpha_m-0.85\)。主表只引用第\ref{ch:result}章。

求解器只在数值核里。页面与附录共用种子、厚度和策略名称；图\ref{fig:demo}--\ref{fig:demothick} 取自同一条状态流的起始帧、探针窗、品质段、开门全貌、对照开环、停机附近以及 50 mm 时域终点。晚期帧对应品质加热、冷点 \(70.21^\circ\mathrm{C}\)，还没有切入保持。

\begin{figure}[htbp]
\centering
\begin{subfigure}[t]{0.48\textwidth}
\centering
\includegraphics[width=\linewidth]{E_demo_start.png}
\caption{起始帧。色带接近初温，蒸汽为基线。}
\label{fig:demo1}
\end{subfigure}
\hfill
\begin{subfigure}[t]{0.48\textwidth}
\centering
\includegraphics[width=\linewidth]{E_demo_probe.png}
\caption{探针窗前半段。蒸汽从喷嘴喷出。}
\label{fig:demo2}
\end{subfigure}

\vspace{0.55em}
\begin{subfigure}[t]{0.72\textwidth}
\centering
\includegraphics[width=\linewidth]{E_demo_late.png}
\caption{1680 s 品质段。冷点 \(70.21^\circ\mathrm{C}\)，进度 0.558。}
\label{fig:demo3}
\end{subfigure}
\caption{组合档、对偶策略、肉糜 30 mm、种子 7。页面读取同一数值核写出的状态流。}
\label{fig:demo}
\end{figure}

对照开环时，底部曲线用虚线叠出菜单切断后的冷点与进度，右栏同时给出对照本帧，见图\ref{fig:democmp}。开环在 1170 s 切断加热，冷点停在 \(16^\circ\mathrm{C}\) 量级；对偶在同一时刻仍走品质加热。事件条把滑条送到探针窗、1680 s 和停机 1866 s，用来核对附录抽样，而不是另跑一条积分。

\begin{figure}[htbp]
\centering
\includegraphics[width=0.92\textwidth]{E_demo_compare.png}
\caption{对照开环菜单。虚线为开环冷点与进度；实线为对偶。组合档、30 mm、1680 s。}
\label{fig:democmp}
\end{figure}

图\ref{fig:demoopen}给出门打开时的腔体：顶部六根 U 形加热管带陶瓷端帽、后壁十叶深色扭曲风机与近黑挡板、十六块几乎贴合的圆角食品薄片（第 5 片为脂肪夹层的更深色带）和底部带双轴标注的曲线。这一帧是开环菜单、组合档、起始时刻。拖动旋转、滚轮缩放；默认方位角 0.52、俯角 0.20、距离 6.55。加热管自发光跟加热条走，探针窗前半段蒸汽粒子从左侧喷嘴喷出，风机叶片转速跟风机条走，排汽百叶随排汽条张开。水位只出现在水箱液面，门只出现在铰链开度。

\begin{figure}[htbp]
\centering
\includegraphics[width=0.92\textwidth]{E_demo_open.png}
\caption{门打开时的腔体。开环菜单、组合档、起始帧。页面读取已经写出的状态流。}
\label{fig:demoopen}
\end{figure}

\begin{figure}[htbp]
\centering
\includegraphics[width=0.92\textwidth]{E_demo_stop.png}
\caption{停机附近。组合档、对偶、30 mm。页面显示 1864 s、\(e_{\alpha}=+0.145\)；核日志停机 1866 s。}
\label{fig:demostop}
\end{figure}

\begin{figure}[htbp]
\centering
\includegraphics[width=0.92\textwidth]{E_demo_thick.png}
\caption{厚度 50 mm、对偶、组合档、时域终点。停机 3298 s，\(e_{\alpha}=-0.782\)，低于设计线；方差已是 0.224。}
\label{fig:demothick}
\end{figure}
\FloatBarrier
